\section{Explicación detallada de la operación de convolución}

\subsection{Estudio de caso: detección de la letra X}

Amini usó un ejemplo didáctico clásico para explicar de manera intuitiva la operación de convolución. Supongamos que queremos determinar si una imagen contiene la letra "X". Aunque dos "X" se vean distintas después de un desplazamiento, un escalamiento o una rotación, la computadora debería poder reconocer ambas como "X".

\begin{figure}[H]
\centering
\includegraphics[width=0.85\textwidth]{slides-images/slide-27.jpg}
\caption{¿X o X? La imagen se representa como una matriz de píxeles, pero la computadora la compara píxel por píxel: incluso para la misma "X", después de un desplazamiento la matriz de valores de píxeles es completamente distinta.\protect\footnotemark}
\end{figure}
\footnotetext{Diapositiva 27. Referencia: Rohrer, "How do CNNs work?".}

Solución: en lugar de comparar a nivel global, se compara a nivel de \textbf{características locales}. Una "X" está formada por varias características: líneas diagonales, puntos de intersección, etcétera. Basta con encontrar estas características locales en ambas imágenes para considerar que ambas son "X".

\begin{figure}[H]
\centering
\includegraphics[width=0.9\textwidth]{slides-images/slide-28.jpg}
\caption{Características de la letra X: segmentos diagonales y un punto de intersección. Aunque las dos X estén ubicadas en posiciones globales distintas, sus características locales coinciden.\protect\footnotemark}
\end{figure}
\footnotetext{Diapositiva 28.}

Estas características locales son los \textbf{filtros} (filters). Cada filtro es una pequeña matriz numérica que representa el patrón que queremos detectar.

\begin{figure}[H]
\centering
\includegraphics[width=0.9\textwidth]{slides-images/slide-29.jpg}
\caption{Tres filtros de 3$\times$3 usados para detectar características de "X": detectan, respectivamente, la diagonal izquierda, la intersección en cruz y la diagonal derecha.\protect\footnotemark}
\end{figure}
\footnotetext{Diapositiva 29.}

\subsection{Definición matemática de la operación de convolución}

Los pasos de la operación de convolución son los siguientes:

\begin{enumerate}
    \item Colocar el filtro en alguna posición de la imagen de entrada
    \item Sobre la región cubierta por el filtro, realizar una \textbf{multiplicación elemento a elemento} (element-wise multiplication)
    \item \textbf{Sumar} todos los productos
    \item Escribir el resultado en la posición correspondiente del \textbf{mapa de características} (feature map)
    \item Deslizar el filtro a la siguiente posición y repetir el proceso anterior
\end{enumerate}

\begin{figure}[H]
\centering
\includegraphics[width=0.9\textwidth]{slides-images/slide-31.jpg}
\caption{Ejemplo de la operación de convolución: convolución de una imagen de entrada de 5$\times$5 con un filtro de 3$\times$3. El filtro se desliza sobre la entrada, se multiplica elemento a elemento y se suma, obteniendo el mapa de características.\protect\footnotemark}
\end{figure}
\footnotetext{Diapositiva 31. Referencia: Karn, "Intuitive CNNs".}

Tomando como ejemplo una imagen de entrada de $5 \times 5$ y un filtro de $3 \times 3$, el proceso de convolución genera un mapa de características de $3 \times 3$:

\begin{figure}[H]
\centering
\includegraphics[width=0.9\textwidth]{slides-images/slide-32.jpg}
\caption{Primer paso de la convolución: el filtro cubre la región de 3$\times$3 de la esquina superior izquierda; al multiplicar elemento a elemento y sumar se obtiene 4, que se escribe en la primera posición del mapa de características.\protect\footnotemark}
\end{figure}
\footnotetext{Diapositiva 32.}

Para el caso general, la operación de convolución puede expresarse con la siguiente fórmula:

$$y_{p,q} = \sum_{i=1}^{k} \sum_{j=1}^{k} w_{i,j} \cdot x_{i+p, j+q} + b$$

donde:
\begin{itemize}
    \item $x$ es la imagen de entrada
    \item $w$ es la matriz de pesos del filtro, de tamaño $k \times k$
    \item $b$ es el término de sesgo
    \item $y_{p,q}$ es el valor de salida del mapa de características en la posición $(p,q)$
    \item $(p,q)$ es la posición de la neurona en la capa oculta
\end{itemize}

\begin{figure}[H]
\centering
\includegraphics[width=0.85\textwidth]{slides-images/slide-40.jpg}
\caption{Resultado completo de la convolución: la convolución de una entrada de 5$\times$5 con un filtro de 3$\times$3 (stride=1) produce un mapa de características de 3$\times$3.\protect\footnotemark}
\end{figure}
\footnotetext{Diapositiva 40.}

\begin{knowledgebox}{La esencia de la convolución es la "suma ponderada"}
Si entendiste la propagación hacia adelante de la capa totalmente conectada de la Clase 1 (realizar una suma ponderada de la entrada y luego sumar un sesgo), entonces lo que hace la capa convolucional es, en esencia, exactamente lo mismo. La única diferencia es que cada neurona de la capa totalmente conectada ve \textbf{toda} la entrada, mientras que cada neurona de la capa convolucional solo ve la entrada dentro de una \textbf{ventana local}, y los pesos de esa ventana se comparten espacialmente.
\end{knowledgebox}

\subsection{Los filtros producen distintos mapas de características}

Filtros con distintos pesos producen mapas de características completamente diferentes. Al cambiar los valores del filtro, se puede lograr:

\begin{figure}[H]
\centering
\includegraphics[width=0.9\textwidth]{slides-images/slide-41.jpg}
\caption{Distintos filtros producen distintos efectos: nitidez (Sharpen), detección de bordes (Edge Detect) y detección de bordes fuerte (Strong Edge Detect).\protect\footnotemark}
\end{figure}
\footnotetext{Diapositiva 41.}

\begin{itemize}
    \item \textbf{Filtro de nitidez}: aumenta el peso del píxel central y disminuye el de los píxeles circundantes, resaltando los detalles de la imagen
    \item \textbf{Filtro de detección de bordes}: es, en esencia, un operador de derivada que detecta los cambios abruptos de brillo de alto a bajo
    \item \textbf{Filtro de detección de bordes fuerte}: un operador de derivada más agresivo, que produce bordes más marcados
\end{itemize}

\begin{knowledgebox}{Del diseño manual al aprendizaje automático de los filtros}
Antes del aprendizaje profundo, durante mucho tiempo se usaron filtros diseñados manualmente (como el operador de Sobel o el operador Laplaciano) para detectar patrones específicos. El avance clave del aprendizaje profundo fue: \textbf{hacer que los pesos del filtro sean parámetros que se pueden aprender}, optimizándolos automáticamente a partir de los datos mediante retropropagación y descenso de gradiente. Así, la red puede descubrir por sí misma las características más útiles para la tarea en cuestión, sin necesidad de diseño humano.
\end{knowledgebox}

\subsection{Cálculo del tamaño de la salida}

El tamaño de salida de la operación de convolución depende de tres hiperparámetros:

\begin{table}[H]
\centering
\begin{tabular}{lp{10cm}}
\toprule
\textbf{Hiperparámetro} & \textbf{Descripción} \\
\midrule
Tamaño del filtro (kernel size) & Las dimensiones espaciales del filtro, por ejemplo $3 \times 3$, $5 \times 5$ \\
Paso (stride) & El número de píxeles que se desplaza el filtro en cada paso. stride=1 significa que se mueve un píxel cada vez; stride=2 significa que salta dos píxeles cada vez \\
Relleno (padding) & Si se agregan ceros en los bordes de la entrada. Si no se rellena, el tamaño de salida se reduce \\
\bottomrule
\end{tabular}
\caption{Hiperparámetros clave de la capa convolucional}
\end{table}

Para una entrada de tamaño $n$, un filtro de tamaño $k$ y un paso de $s$, el tamaño de salida es:

$$\text{output size} = \left\lfloor \frac{n - k}{s} \right\rfloor + 1$$

\subsection{Resumen de la sección}

La operación de convolución desliza un filtro sobre la imagen, realiza multiplicación elemento a elemento y suma, generando un mapa de características. Distintos filtros detectan distintos patrones. En el aprendizaje profundo, los pesos de los filtros se optimizan automáticamente mediante el entrenamiento, en lugar de diseñarse manualmente. Las tres ventajas centrales de la convolución son: conectividad local, uso compartido de parámetros y preservación de la estructura espacial.

\newpage
